% ============================================================
% introduction.tex
% ============================================================

\section{Introduction}
\label{sec:introduction}

The human cerebral cortex undergoes measurable structural change throughout the adult lifespan. Cortical thickness decreases in a regionally heterogeneous pattern that accelerates after the fifth decade~\cite{fjell2010structural, salat2004thinning}, cortical surface area follows a largely independent developmental trajectory~\cite{storsve2014differential}, and cortical volume, which reflects the product of the two, declines across nearly every parcellated region by late adulthood~\cite{raz2005regional}. These morphometric indices are not merely markers of biological aging; they covary with cognitive ability~\cite{deary2010neuroscience}, affective temperament~\cite{riccelli2017surface}, and health-relevant behaviors such as habitual sleep quality~\cite{sexton2014poor}. A growing body of work therefore treats cortical morphometry as a common endpoint through which disparate behavioral and demographic factors can be compared on a shared anatomical scale.

Despite this convergence of interest, most studies examine a single behavioral predictor in a single sample and report effects in one cortical metric. Sleep quality has been linked to regional thinning in temporal and frontal areas in community-dwelling older adults~\cite{sexton2014poor, branger2016relationships}, but few studies have tested the same association in a young-adult cohort with restricted age variance, where the contribution of age-related confounds is minimized. Fluid intelligence correlates positively with total brain volume and, to a lesser extent, with regional gray matter density~\cite{pietschnig2015meta, colom2009human}; yet the specific relationship between intelligence and cortical \textit{surface area}, as opposed to thickness, remains underexplored in large samples despite theoretical grounds for expecting such an association~\cite{jung2007pfit}. Neuroticism, the personality dimension most consistently linked to negative affect and stress reactivity~\cite{costa1992neo}, has been associated with both thicker and thinner cortices depending on participant age, brain region, and imaging modality~\cite{bjornebekk2013neuronal, jackson2011cortical}, leaving open the question of whether these effects are present in young adulthood or instead emerge only with prolonged exposure to stress-related physiological processes~\cite{lupien2009effects}. In short, there is no single study that examines sleep, cognition, and personality within the same analytic framework while simultaneously characterizing age-related cortical decline across a complementary lifespan sample.

The present work addresses this gap through a dual-cohort design. We use the Human Connectome Project Young Adult (HCP-YA) dataset~\cite{vanessen2013wuminn}, which provides high-resolution structural MRI and extensive behavioral phenotyping for 1{,}104 adults aged 22 to 37, to investigate three research questions: whether subjective sleep quality (Pittsburgh Sleep Quality Index; PSQI~\cite{buysse1989pittsburgh}) predicts regional cortical thickness, whether fluid intelligence (Penn Progressive Matrices) predicts cortical surface area, and whether dispositional neuroticism (NEO Five-Factor Inventory~\cite{costa1992neo}) predicts either metric. For the HCP cohort, structural data are parcellated using the Desikan-Killiany atlas~\cite{desikan2006automated}, yielding 34 regions per hemisphere (68 total). We then turn to the Aging Brain and Behaviour Cohort (AABC Release~2)~\cite{bookheimer2019aging}, which includes 1{,}150 adults spanning ages 36 to 89, parcellated with the finer-grained Glasser HCP-MMP1 atlas~\cite{glasser2016multimodal} (180 regions per hemisphere). In this older cohort, we characterize the regional pattern of age-related thickness and volume decline and test whether sleep quality and neuroticism continue to predict cortical structure after the dominant contribution of age has been accounted for.

Across both cohorts, we apply a uniform analytic strategy. Group comparisons use Welch's \textit{t}-test, continuous associations are assessed with Pearson's correlation, and all \textit{p}-values are adjusted for multiplicity using the Benjamini-Hochberg false discovery rate procedure at \textit{q}\,<\,.05~\cite{benjamini1995controlling}. Effect sizes are reported as Cohen's \textit{d} for group contrasts and as Pearson's \textit{r} for correlational analyses, with 95\% confidence intervals throughout. This consistent framework enables direct comparison of effect magnitudes across research questions.

Three principal findings emerge from this work. First, fluid intelligence showed a positive association with cortical surface area in every region tested, with the largest effects concentrated in the inferior and middle temporal cortex. Second, cortical thickness and volume declined with age across nearly all 360 Glasser regions, with primary motor cortex exhibiting the single strongest effect (\textit{r}\,=\,$-$0.69 for thickness). Third, neuroticism was unrelated to cortical thickness in young adults (0 of 68 regions survived FDR correction) but was associated with widespread thinning in the aging cohort (327 of 360 regions, Cohen's \textit{d} up to $-$0.39), suggesting that the structural correlates of this personality trait accumulate with age. Conversely, sleep quality showed a modest three-region effect in young adults and no surviving associations in the aging sample, where age itself accounted for the overwhelming majority of variance in cortical structure.

These findings carry two implications for the broader field. Methodologically, they demonstrate that the same behavioral variable can yield opposing conclusions depending on the age composition of the sample and the cortical metric under examination, reinforcing the need for multi-cohort, multi-metric designs. Substantively, the divergent age profiles for neuroticism and sleep quality suggest that the biological pathways linking personality to cortical structure may involve chronic, cumulative processes, such as hypothalamic-pituitary-adrenal axis dysregulation and elevated allostatic load~\cite{mcewen2006protective}, that require decades to leave a detectable morphometric signature.
